% CRITICAL DISCLAIMER: This draft is a structural aid only. The author must
% contribute original research. The AI assists with organization and phrasing
% but cannot replace genuine scholarly contribution. The author is responsible
% for verifying correctness, novelty, and ethical standards.

\documentclass[twoside,11pt]{article}

% Any additional packages needed should be included after jmlr2e.
% Note that jmlr2e.sty includes epsfig, amssymb, natbib and graphicx,
% and defines many common macros, such as 'proof' and 'example'.
%
% It also sets the bibliographystyle to plainnat; for more information on
% natbib citation styles, see the natbib documentation.

\usepackage{jmlr2e}
\usepackage{amsmath}
\usepackage{amsfonts}
\usepackage{amsthm}
\usepackage{algorithm}
\usepackage{algorithmic}
\usepackage{booktabs}
\usepackage{multirow}
\usepackage{url}

\hypersetup{hidelinks}

% ---------------------------------------------------------------
% Theorem environments
% ---------------------------------------------------------------
\newtheorem{theorem}{Theorem}[section]
\newtheorem{lemma}[theorem]{Lemma}
\newtheorem{proposition}[theorem]{Proposition}
\newtheorem{corollary}[theorem]{Corollary}
\newtheorem{definition}[theorem]{Definition}
\newtheorem{remark}[theorem]{Remark}

% ---------------------------------------------------------------
% Custom macros
% ---------------------------------------------------------------
\newcommand{\Vocab}{\Sigma}
\newcommand{\LM}{\mathcal{P}}
\newcommand{\Oracle}{\mathcal{O}}
\newcommand{\Constraint}{\mathcal{C}}
\newcommand{\Logic}{\mathcal{L}}
\newcommand{\KL}{\mathrm{KL}}
\newcommand{\TV}{\mathrm{TV}}
\newcommand{\E}{\mathbb{E}}
\newcommand{\R}{\mathbb{R}}
\newcommand{\N}{\mathbb{N}}
\newcommand{\indicator}{\mathbf{1}}
\newcommand{\support}{\mathrm{supp}}
\newcommand{\SFS}{\mathrm{SFS}}
\newcommand{\EFG}{\mathrm{EFG}}

% ---------------------------------------------------------------
% JMLR heading (placeholder values — editor fills in final version)
% ---------------------------------------------------------------
\usepackage{lastpage}
\jmlrheading{1}{2026}{1-\pageref{LastPage}}{4/26}{10/26}{26-0001}{[Author Names]}

% Short headings: running head ≤50 characters
\ShortHeadings{Semantically Constrained Decoding}{[Author Names]}
\firstpageno{1}

% ---------------------------------------------------------------
\begin{document}

\title{Semantically Constrained Decoding: A Formal Theory of\\
Distribution-Aligned Neurosymbolic Generation}

% ANONYMIZATION: Remove ALL acknowledgments, affiliations, and self-cites
% that reveal identity for review. Add these only in camera-ready.
\author{\name [Author Name] \email [email]\\
       \addr [Affiliation]\\
       [Address]}

\editor{[Action Editor]}

\maketitle

% ---------------------------------------------------------------
% ABSTRACT
% ---------------------------------------------------------------
\begin{abstract}%
Constrained decoding has emerged as the dominant paradigm for enforcing
structural correctness in large language model outputs. However, all
existing theoretical analyses of constrained decoding---including
distribution-alignment guarantees and convergence proofs---are restricted
to syntactic constraints such as context-free grammars and regular
expressions. We extend constrained decoding theory to semantic
constraints defined by satisfiability modulo theories (SMT) and
first-order logic (FOL) oracles. We formalize semantically constrained
decoding as sampling from the language model's true conditional
distribution given a semantic constraint oracle, and prove that naive
semantic filtering---the direct analogue of grammar-constrained
decoding---introduces distribution distortion whose magnitude depends on
the constraint's prefix-decidability. We propose SCD-MH, a
Metropolis-Hastings-based sampling algorithm that provably converges to
the true conditional distribution under mild conditions, and
characterize its convergence rate as a function of constraint complexity
and model quality. We further generalize recent results on reasoning
preservation under constrained generation from syntactic to semantic
settings. This work provides the first unified formal theory of
distribution-aligned constrained decoding spanning both syntactic and
semantic constraints.
\end{abstract}

\begin{keywords}
constrained decoding, neurosymbolic AI, formal verification, distribution alignment, Metropolis-Hastings sampling
\end{keywords}

% ===============================================================
% 1. INTRODUCTION
% ===============================================================
\section{Introduction}
\label{sec:introduction}

Large language models (LLMs) have demonstrated remarkable fluency across a
wide range of generation tasks, from open-ended dialogue to code synthesis
and mathematical reasoning \citep{brown2020language, touvron2023llama}.
Despite this fluency, LLMs frequently produce outputs that are logically
incorrect, type-inconsistent, or violate domain-specific constraints
\citep{chen2021evaluating, mirzadeh2024gsm}. This tension between
fluency and correctness has motivated a growing body of work on
\emph{constrained decoding}: modifying the generation process at inference
time to ensure that outputs satisfy specified structural or logical
requirements.

The most mature line of constrained decoding research addresses
\emph{syntactic constraints}---requirements expressible as membership in a
formal language defined by a context-free grammar (CFG) or regular
expression. Grammar-constrained decoding (GCD) enforces such constraints by
masking logits at each generation step to exclude tokens that would render
the partial output unextendable to a string in the target language
\citep{hokamp2017lexically, post2018fast, scholak2021picard}. This
approach has been widely adopted in structured generation tasks, including
JSON output formatting, SQL synthesis, and code generation
\citep{beurerkellner2024domino}.

A pivotal theoretical advance was the work of \citet{ugare2024gad}, who
demonstrated that GCD does \emph{not} sample from the LLM's true conditional
distribution $P(\cdot \mid x \in L(G))$, where $L(G)$ is the language of the
constraining grammar. Instead, GCD samples from a \emph{distorted}
distribution arising from the path-dependent interaction between per-step
masking and autoregressive factorization. They proposed the ASAp (Adaptive
Sampling with Approximate expected futures) algorithm, which uses importance
weights derived from ``expected future grammaticality'' to correct this
distortion, providing the first convergence guarantees for syntactically
constrained decoding. Concurrently, \citet{banerjee2025crane} provided a
theoretical explanation for the empirical observation that restrictive
grammars can degrade reasoning performance: constraining the output space
eliminates intermediate tokens needed for chain-of-thought reasoning, and
they showed that grammar augmentation can preserve reasoning capacity.

These advances represent significant progress. However, a critical gap
remains: \textbf{all existing constrained decoding theory is restricted to
syntactic constraints}. In practice, many of the most important correctness
requirements are \emph{semantic} in nature---they concern the meaning or
logical validity of the output, not merely its syntactic form. Examples
include:
\begin{itemize}
    \item Logical consistency of natural language arguments, verifiable via
    first-order logic (FOL) theorem provers \citep{olausson2023linc,
    ganguly2024pot};
    \item Arithmetic and algebraic correctness of mathematical reasoning,
    checkable via SMT solvers \citep{demoura2008z3};
    \item Type correctness and runtime safety of generated code, enforceable
    via type checkers and symbolic execution \citep{chen2021evaluating};
    \item Relational consistency in structured knowledge extraction,
    verifiable via Prolog or Datalog engines \citep{ma2025logically}.
\end{itemize}

A number of practical neurosymbolic systems have been proposed to enforce
such constraints, including LINC \citep{olausson2023linc}, Proof of Thought
\citep{ganguly2024pot}, VERGE \citep{verge2026}, and Logically Constrained
Decoding \citep{ma2025logically}. These systems achieve impressive empirical
results by coupling LLM generation with symbolic verification. However,
\emph{none of them provides formal guarantees on distribution alignment or
convergence}. We lack answers to fundamental questions: Does naive semantic
filtering distort the LLM's distribution? If so, by how much? Can we
correct the distortion while maintaining the constraint? How does constraint
complexity affect convergence?

\paragraph{Contributions.} In this paper, we provide the first formal theory
of semantically constrained decoding. Our contributions are:
\begin{itemize}
    \item We \textbf{formalize semantically constrained decoding} as
    sampling from an LLM's conditional distribution given a semantic
    constraint oracle, introducing the \emph{prefix-decidability spectrum}
    that characterizes the fundamental difference between syntactic and
    semantic constraints (Section~\ref{sec:formulation}).

    \item We \textbf{prove that naive semantic filtering introduces
    distribution distortion} analogous to GCD's distortion, and
    characterize the distortion gap as a function of \emph{semantic future
    satisfiability}---a generalization of the ``expected future
    grammaticality'' concept from \citet{ugare2024gad}
    (Section~\ref{sec:distortion}).

    \item We \textbf{propose SCD-MH} (Semantically Constrained Decoding
    via Metropolis-Hastings), a sampling algorithm with provable convergence
    to the true conditional distribution, and characterize its mixing time
    as a function of constraint complexity and LLM quality
    (Section~\ref{sec:scdmh}).

    \item We \textbf{generalize reasoning preservation results} from
    syntactic to semantic constraints, showing that for any prefix-decidable
    semantic constraint, there exists an augmented constraint that preserves
    reasoning capacity while ensuring semantic correctness
    (Section~\ref{sec:reasoning}).

    \item We \textbf{validate our theory empirically} on logical reasoning,
    arithmetic, multi-step deduction, and typed code generation benchmarks,
    demonstrating that SCD-MH corrects distribution distortion and improves
    task accuracy over naive filtering (Section~\ref{sec:experiments}).
\end{itemize}

% ===============================================================
% 2. PRELIMINARIES
% ===============================================================
\section{Preliminaries}
\label{sec:preliminaries}

\subsection{Language Models and Autoregressive Generation}
\label{sec:prelim-lm}

Let $\Vocab$ denote a finite vocabulary of tokens, and let $\Vocab^*$ denote
the set of all finite sequences over $\Vocab$, including the empty sequence
$\epsilon$. We use $\texttt{EOS} \in \Vocab$ to denote the distinguished
end-of-sequence token. An autoregressive language model defines a probability
distribution $P$ over $\Vocab^*$ via the chain rule factorization:
\begin{equation}
\label{eq:autoregressive}
P(x_{1:n}) = \prod_{t=1}^{n} P(x_t \mid x_{1:t-1}),
\end{equation}
where $x_{1:0} = \epsilon$ and we assume $x_n = \texttt{EOS}$ for all
complete sequences. We write $P(x_t \mid x_{1:t-1})$ for the conditional
next-token distribution at step $t$.

\subsection{Constrained Decoding}
\label{sec:prelim-cd}

A \emph{constraint} is a subset $\Constraint \subseteq \Vocab^*$ of
acceptable sequences. Given a language model $P$ and a constraint
$\Constraint$, the goal of constrained decoding is to sample from the
\emph{true conditional distribution}:
\begin{equation}
\label{eq:true-conditional}
P^*(x) \triangleq P(x \mid x \in \Constraint) = \frac{P(x) \cdot
\indicator[x \in \Constraint]}{\sum_{x' \in \Constraint} P(x')},
\end{equation}
assuming $P(\Constraint) \triangleq \sum_{x' \in \Constraint} P(x') > 0$.

The dominant approach to enforcing constraints during autoregressive
generation is \emph{token masking}. At each step $t$, given prefix
$x_{1:t-1}$, the set of \emph{allowed tokens} is defined as:
\begin{equation}
\label{eq:masking-operator}
A_\Constraint(x_{1:t-1}) = \{v \in \Vocab : \exists\, y \in \Vocab^*
\text{ s.t. } x_{1:t-1} \cdot v \cdot y \in \Constraint\},
\end{equation}
and the masked distribution at step $t$ is:
\begin{equation}
\label{eq:masked-dist}
Q(x_t \mid x_{1:t-1}) = \frac{P(x_t \mid x_{1:t-1}) \cdot
\indicator[x_t \in A_\Constraint(x_{1:t-1})]}{\sum_{v \in
A_\Constraint(x_{1:t-1})} P(v \mid x_{1:t-1})}.
\end{equation}
Grammar-constrained decoding (GCD) implements this procedure when
$\Constraint = L(G)$ for a context-free grammar $G$, using the grammar's
incremental parsing state to compute $A_{L(G)}$ efficiently.

\subsection{Grammar-Aligned Decoding}
\label{sec:prelim-gad}

\citet{ugare2024gad} proved that sampling from the step-wise masked
distribution $Q$ (Equation~\ref{eq:masked-dist}) does \emph{not} yield
samples from the true conditional $P^*$ (Equation~\ref{eq:true-conditional}).
The key insight is that the masking operator renormalizes at each step
independently, failing to account for the \emph{future} probability of
completing a prefix to a valid sequence. Formally, define the \emph{expected
future grammaticality} of a prefix $x_{1:t}$ as:
\begin{equation}
\label{eq:efg}
\EFG(x_{1:t}) = P(\exists\, y \in \Vocab^* : x_{1:t} \cdot y \in
\Constraint \mid x_{1:t}),
\end{equation}
which measures the probability, under the unconstrained model, that the
prefix $x_{1:t}$ can be completed to a sequence in $\Constraint$. The
distribution $Q$ induced by step-wise masking satisfies:
\begin{equation}
\label{eq:gcd-distortion}
Q(x_{1:n}) = P^*(x_{1:n}) \cdot \prod_{t=1}^{n}
\frac{\EFG(x_{1:t-1})}{\EFG(x_{1:t})},
\end{equation}
which equals $P^*$ only when the EFG ratio is constant across all
paths---a condition that is generically violated. The ASAp algorithm
proposed by \citet{ugare2024gad} corrects this distortion by incorporating
approximate EFG estimates as importance weights.

\subsection{Symbolic Constraint Oracles}
\label{sec:prelim-oracle}

We generalize beyond grammar-based constraints by introducing
\emph{symbolic constraint oracles}. Let $\Logic$ be a formal logic (e.g.,
propositional logic, first-order logic, or SMT-LIB). A \emph{semantic
constraint} is specified by a formula $\varphi \in \Logic$ together with
an interpretation function $\mathcal{I}$ that maps token sequences to
logical structures.

\begin{definition}[Symbolic Constraint Oracle]
\label{def:oracle}
A symbolic constraint oracle for a formula $\varphi \in \Logic$ is a
function $\Oracle_\varphi : \Vocab^* \to \{\textsc{Sat}, \textsc{Unsat},
\textsc{Unknown}\}$ such that:
\begin{enumerate}
    \item If $\Oracle_\varphi(x) = \textsc{Sat}$, then $\mathcal{I}(x)
    \models \varphi$ (soundness).
    \item If $\Oracle_\varphi(x) = \textsc{Unsat}$, then $\mathcal{I}(x)
    \not\models \varphi$ (soundness).
    \item For complete sequences, $\Oracle_\varphi(x) \in \{\textsc{Sat},
    \textsc{Unsat}\}$ (completeness on termination).
\end{enumerate}
\end{definition}

The key distinction from syntactic oracles (grammar membership tests) is
that semantic oracles may return $\textsc{Unknown}$ on \emph{prefixes}: it
may be undecidable or computationally infeasible to determine whether a
partial output can be extended to satisfy $\varphi$. This is the fundamental
source of the additional complexity in semantically constrained decoding.

Concrete instantiations include Z3 \citep{demoura2008z3} for SMT
constraints, Prolog engines for relational constraints, and FOL theorem
provers for logical validity constraints \citep{olausson2023linc}.

% ===============================================================
% 3. PROBLEM FORMULATION
% ===============================================================
\section{Semantically Constrained Decoding: Problem Formulation}
\label{sec:formulation}

\subsection{Semantic Constraints as Measurable Sets}
\label{sec:measurable}

Given a semantic constraint specified by a formula $\varphi \in \Logic$ and
an oracle $\Oracle_\varphi$, we define the \emph{semantic constraint set}:
\begin{equation}
\label{eq:semantic-constraint-set}
\Constraint_\varphi = \{x \in \Vocab^* : \Oracle_\varphi(x) = \textsc{Sat}\}.
\end{equation}
We equip $\Vocab^*$ with the discrete $\sigma$-algebra and the probability
measure induced by $P$. Under this measure, $\Constraint_\varphi$ is
trivially measurable, and we assume throughout that
$P(\Constraint_\varphi) > 0$ (the LLM assigns positive probability to at
least one satisfying sequence).

Unlike syntactic constraint sets $L(G)$, which admit efficient incremental
membership tests via parsing automata, the set $\Constraint_\varphi$ may
not admit efficient prefix-based membership tests. This motivates the
following classification.

\subsection{The Prefix-Decidability Spectrum}
\label{sec:prefix-decidability}

The efficiency and correctness of constrained decoding depend critically on
the extent to which the constraint can be evaluated on \emph{partial}
outputs. We introduce a hierarchy that captures this.

\begin{definition}[Prefix Decidability]
\label{def:prefix-decidability}
A semantic constraint $\Constraint_\varphi$ is:
\begin{enumerate}
    \item \textbf{Prefix-decidable} if for all prefixes $x_{1:t}$, the
    oracle can determine in finite time whether there exists a completion
    $y$ such that $x_{1:t} \cdot y \in \Constraint_\varphi$:
    \begin{equation}
    \label{eq:prefix-decidable}
    \forall\, x_{1:t} \in \Vocab^* : \Oracle_\varphi^{\mathrm{prefix}}(x_{1:t})
    \in \{\textsc{Extendable}, \textsc{Dead}\},
    \end{equation}
    where $\textsc{Extendable}$ means $\exists\, y : x_{1:t} \cdot y \in
    \Constraint_\varphi$ and $\textsc{Dead}$ means no such $y$ exists.

    \item \textbf{Prefix-approximable} if $\Oracle_\varphi^{\mathrm{prefix}}$
    can return a sound but incomplete answer: $\textsc{Dead}$ implies no
    extension exists, but $\textsc{Extendable}$ or $\textsc{Unknown}$ may
    be returned when the true answer is $\textsc{Dead}$.

    \item \textbf{Prefix-opaque} if no efficient prefix oracle exists; the
    constraint can only be evaluated on complete sequences.
\end{enumerate}
\end{definition}

Syntactic constraints (CFGs, regular expressions) are prefix-decidable by
construction: the parsing automaton provides an exact prefix oracle. Many
semantic constraints are only prefix-approximable---for instance, checking
whether a partial arithmetic expression can be extended to a true equation
requires reasoning about ungenerated terms. Some constraints are
prefix-opaque---for example, verifying that a complete proof is valid
requires the full proof text.

\subsection{The Semantic Constrained Decoding Problem}
\label{sec:scd-problem}

We can now state the central problem addressed in this paper.

\begin{definition}[Semantic Constrained Decoding Problem]
\label{def:scd}
Given an autoregressive language model $P$ over $\Vocab^*$, a semantic
constraint $\varphi \in \Logic$ with oracle $\Oracle_\varphi$, and the
induced constraint set $\Constraint_\varphi$, the \emph{semantic constrained
decoding problem} is to sample from the target distribution:
\begin{equation}
\label{eq:target-distribution}
P^*_\varphi(x) \triangleq P(x \mid x \in \Constraint_\varphi) =
\frac{P(x) \cdot \indicator[x \in \Constraint_\varphi]}{P(\Constraint_\varphi)}.
\end{equation}
\end{definition}

This generalizes the grammar-aligned decoding problem of
\citet{ugare2024gad} from $\Constraint = L(G)$ to $\Constraint =
\Constraint_\varphi$ for arbitrary semantic constraints. The challenge is
that the normalizing constant $P(\Constraint_\varphi)$ is generally
intractable, and the prefix-based filtering techniques that work for
syntactic constraints may not apply.

% ===============================================================
% 4. DISTRIBUTION DISTORTION
% ===============================================================
\section{Distribution Distortion Under Semantic Constraints}
\label{sec:distortion}

This section contains our core theoretical results. We show that naive
semantic filtering---the direct generalization of GCD to semantic
constraints---introduces distribution distortion, characterize its
magnitude, and identify conditions under which it vanishes.

\subsection{Naive Semantic Filtering}
\label{sec:naive-filtering}

The most natural approach to enforcing a semantic constraint during
autoregressive generation is to apply the masking operator
(Equation~\ref{eq:masking-operator}) using the semantic oracle. At each step
$t$, given prefix $x_{1:t-1}$, define:
\begin{equation}
\label{eq:semantic-allowed}
A_\varphi(x_{1:t-1}) = \{v \in \Vocab :
\Oracle_\varphi^{\mathrm{prefix}}(x_{1:t-1} \cdot v) \neq \textsc{Dead}\},
\end{equation}
and sample from the renormalized distribution:
\begin{equation}
\label{eq:naive-semantic-filter}
Q_\varphi(x_t \mid x_{1:t-1}) = \frac{P(x_t \mid x_{1:t-1}) \cdot
\indicator[x_t \in A_\varphi(x_{1:t-1})]}{\sum_{v \in A_\varphi(x_{1:t-1})}
P(v \mid x_{1:t-1})}.
\end{equation}
This is the semantic analogue of GCD. We call it \emph{naive semantic
filtering} (NSF). For prefix-opaque constraints where no prefix oracle is
available, NSF reduces to unconstrained generation followed by
rejection---accepting only complete sequences in $\Constraint_\varphi$.

\subsection{Distortion Theorem}
\label{sec:distortion-theorem}

We now prove that NSF does not sample from the target distribution
$P^*_\varphi$.

\begin{definition}[Semantic Future Satisfiability]
\label{def:sfs}
The \emph{semantic future satisfiability} of a prefix $x_{1:t}$ with respect
to constraint $\varphi$ is:
\begin{equation}
\label{eq:sfs}
\SFS_\varphi(x_{1:t}) = P(\exists\, y \in \Vocab^* : x_{1:t} \cdot y \in
\Constraint_\varphi \mid x_{1:t}).
\end{equation}
\end{definition}

This generalizes the expected future grammaticality (Equation~\ref{eq:efg})
from syntactic to semantic constraints.

\begin{theorem}[Distortion Under Naive Semantic Filtering]
\label{thm:distortion}
Let $P$ be an autoregressive language model, $\varphi$ a semantic constraint
with prefix-decidable or prefix-approximable oracle $\Oracle_\varphi$, and
$Q_\varphi$ the distribution induced by naive semantic filtering
(Equation~\ref{eq:naive-semantic-filter}). Then $Q_\varphi$ satisfies:
\begin{equation}
\label{eq:distorted-distribution}
Q_\varphi(x_{1:n}) = P^*_\varphi(x_{1:n}) \cdot
\prod_{t=1}^{n} \frac{\SFS_\varphi(x_{1:t-1})}{\SFS_\varphi(x_{1:t})},
\end{equation}
for all $x_{1:n} \in \Constraint_\varphi$. Consequently, $Q_\varphi =
P^*_\varphi$ if and only if the ratio
$\SFS_\varphi(x_{1:t-1}) / \SFS_\varphi(x_{1:t})$ is constant across all
prefixes and time steps.
\end{theorem}

\begin{proof}[Proof sketch]
The proof follows the same telescoping argument as the GCD distortion proof
in \citet{ugare2024gad}, generalized to semantic constraints. By
construction, the NSF distribution factorizes as:
\[
Q_\varphi(x_{1:n}) = \prod_{t=1}^{n} Q_\varphi(x_t \mid x_{1:t-1})
= \prod_{t=1}^{n} \frac{P(x_t \mid x_{1:t-1}) \cdot \indicator[x_t
\in A_\varphi(x_{1:t-1})]}{Z_t(x_{1:t-1})},
\]
where $Z_t(x_{1:t-1}) = \sum_{v \in A_\varphi(x_{1:t-1})} P(v \mid
x_{1:t-1})$. By the definition of $\SFS_\varphi$, we have $Z_t(x_{1:t-1})
= \SFS_\varphi(x_{1:t-1}) / \SFS_\varphi(x_{1:t})$ (up to boundary terms
that telescope). Comparing with $P^*_\varphi(x_{1:n}) = P(x_{1:n}) /
P(\Constraint_\varphi)$ yields the stated formula. The product of SFS ratios
is generically non-constant because the probability of future satisfiability
depends on the specific tokens generated, creating path-dependent biases.
% TODO: Complete formal proof with explicit handling of boundary conditions
% and the relationship between Z_t and SFS ratios.
\end{proof}

\subsection{Characterizing the Distortion Gap}
\label{sec:distortion-gap}

Having established that distortion occurs, we now bound its magnitude.

\begin{theorem}[Distortion Bound via Semantic Future Satisfiability]
\label{thm:distortion-bound}
Under the same conditions as Theorem~\ref{thm:distortion}, the KL
divergence between the NSF distribution $Q_\varphi$ and the true conditional
$P^*_\varphi$ is bounded by:
\begin{equation}
\label{eq:kl-bound}
\KL(Q_\varphi \| P^*_\varphi) = \sum_{t=1}^{n}
\E_{x_{1:t} \sim Q_\varphi}\!\left[\log
\frac{\SFS_\varphi(x_{1:t-1})}{\SFS_\varphi(x_{1:t})} \right],
\end{equation}
which equals zero if and only if $\SFS_\varphi(x_{1:t})$ is a deterministic
function of $t$ (not of the specific prefix).
\end{theorem}

\begin{proof}[Proof sketch]
The KL divergence decomposes as:
\[
\KL(Q_\varphi \| P^*_\varphi) = \E_{x \sim Q_\varphi}\!\left[\log
\frac{Q_\varphi(x)}{P^*_\varphi(x)}\right]
= \E_{x \sim Q_\varphi}\!\left[\sum_{t=1}^{n} \log
\frac{\SFS_\varphi(x_{1:t-1})}{\SFS_\varphi(x_{1:t})}\right],
\]
by substituting Equation~\ref{eq:distorted-distribution}. The sum
telescopes partially, but because the expectation is under $Q_\varphi$
(not $P^*_\varphi$), the telescoping does not simplify to a single
term. The bound equals zero iff the SFS ratio is path-independent, which
occurs only when the semantic constraint decomposes into independent per-step
constraints---a condition almost never satisfied by semantic constraints.
% TODO: Complete formal proof with explicit variance decomposition of
% the SFS ratio and connection to mutual information between prefix
% choices and future satisfiability.
\end{proof}

\subsection{When Does Distortion Vanish?}
\label{sec:distortion-vanish}

\begin{proposition}[Distortion Under Extreme Prefix Decidability]
\label{prop:distortion-vanish}
Consider the following cases:
\begin{enumerate}
    \item \textbf{Fully prefix-decidable with independent steps.} If
    $\Constraint_\varphi$ decomposes as $\Constraint_\varphi =
    \Constraint_1 \times \Constraint_2 \times \cdots \times \Constraint_n$
    where $\Constraint_t \subseteq \Vocab$ constrains only the $t$-th
    token independently, then $\KL(Q_\varphi \| P^*_\varphi) = 0$.

    \item \textbf{Terminal-only evaluation.} If $\varphi$ is prefix-opaque
    and NSF reduces to rejection sampling (generate from $P$, accept if
    $x \in \Constraint_\varphi$), then the accepted samples are drawn
    exactly from $P^*_\varphi$ and $\KL(Q_\varphi \| P^*_\varphi) = 0$,
    but the expected number of rejections is $1/P(\Constraint_\varphi)$,
    which may be exponentially large.

    \item \textbf{Prefix-approximable with loose approximation.} If the
    prefix oracle is highly imprecise (many prefixes classified as
    $\textsc{Extendable}$ that are actually $\textsc{Dead}$), the
    distortion approaches zero but computational waste increases.
    Conversely, a tight but imperfect prefix oracle introduces maximal
    distortion.
\end{enumerate}
\end{proposition}

\begin{proof}[Proof sketch]
Case 1 follows because per-step independence makes $\SFS_\varphi(x_{1:t})$
depend only on $t$, not on the specific prefix. Case 2 follows from the
correctness of rejection sampling. Case 3 follows because a loose prefix
oracle makes $A_\varphi(x_{1:t-1})$ close to $\Vocab$, so
$Q_\varphi \approx P$ restricted to $\Constraint_\varphi$ at
termination---the masking has minimal effect. A tight oracle aggressively
prunes the token space, magnifying the path-dependent bias.
% TODO: Complete formal proof with quantitative bounds for Case 3.
\end{proof}

% ===============================================================
% 5. SCD-MH ALGORITHM
% ===============================================================
\section{SCD-MH: A Distribution-Aligned Semantic Decoding Algorithm}
\label{sec:scdmh}

\subsection{Algorithm Design}
\label{sec:algorithm-design}

The distortion results of Section~\ref{sec:distortion} motivate the design
of a corrected sampling algorithm. We propose \emph{SCD-MH} (Semantically
Constrained Decoding via Metropolis-Hastings), which uses the naive
semantically filtered distribution $Q_\varphi$ as a \emph{proposal
distribution} within a Metropolis-Hastings (MH) framework
\citep{metropolis1953equation, hastings1970monte} to converge to the true
conditional $P^*_\varphi$.

The key idea is to generate candidate sequences from $Q_\varphi$ and
accept or reject them based on the MH acceptance criterion computed
using the ratio $P^*_\varphi / Q_\varphi$. Since $P^*_\varphi(x) \propto
P(x) \cdot \indicator[x \in \Constraint_\varphi]$ and $Q_\varphi$ can
be computed exactly (Equation~\ref{eq:naive-semantic-filter}), the
acceptance ratio is tractable.

\begin{algorithm}[t]
\caption{SCD-MH: Semantically Constrained Decoding via Metropolis-Hastings}
\label{alg:scdmh}
\begin{algorithmic}[1]
\REQUIRE Language model $P$, semantic oracle $\Oracle_\varphi$, number of
iterations $T$, burn-in $B$
\ENSURE Samples approximately from $P^*_\varphi$
\STATE Generate initial sequence $x^{(0)} \sim Q_\varphi$ (using naive
semantic filtering; retry if no valid sequence found)
\FOR{$i = 1, \ldots, T$}
    \STATE Generate proposal $x' \sim Q_\varphi$ (full sequence via
    naive semantic filtering)
    \STATE Verify $x' \in \Constraint_\varphi$ using $\Oracle_\varphi$
    \IF{$\Oracle_\varphi(x') = \textsc{Sat}$}
        \STATE Compute acceptance ratio:
        \[
        \alpha(x^{(i-1)}, x') = \min\!\left(1,\;
        \frac{P(x') \cdot Q_\varphi(x^{(i-1)})}{P(x^{(i-1)}) \cdot
        Q_\varphi(x')}\right)
        \]
        \STATE Accept $x^{(i)} \leftarrow x'$ with probability
        $\alpha(x^{(i-1)}, x')$; else $x^{(i)} \leftarrow x^{(i-1)}$
    \ELSE
        \STATE Reject: $x^{(i)} \leftarrow x^{(i-1)}$
    \ENDIF
\ENDFOR
\RETURN $x^{(T)}$ (or $\{x^{(i)}\}_{i > B}$ for multiple samples)
\end{algorithmic}
\end{algorithm}

\begin{remark}
\label{rem:acceptance-ratio}
The acceptance ratio in Algorithm~\ref{alg:scdmh} (line 6) simplifies via
Equation~\ref{eq:distorted-distribution}:
\begin{equation}
\label{eq:mh-acceptance}
\alpha(x, x') = \min\!\left(1,\;
\frac{P(x')}{P(x)} \cdot \frac{Q_\varphi(x)}{Q_\varphi(x')}\right)
= \min\!\left(1,\; \prod_{t=1}^{n}
\frac{\SFS_\varphi(x'_{1:t})}{\SFS_\varphi(x_{1:t})} \cdot
\frac{\SFS_\varphi(x_{1:t-1})}{\SFS_\varphi(x'_{1:t-1})}\right).
\end{equation}
In practice, $P(x)$ and $Q_\varphi(x)$ can be computed exactly from the
model's logits, so the acceptance ratio does not require evaluating
$\SFS_\varphi$ explicitly.
\end{remark}

\subsection{Convergence Guarantee}
\label{sec:convergence}

\begin{theorem}[Convergence of SCD-MH]
\label{thm:convergence}
Let $P$ be an autoregressive language model with $P(\Constraint_\varphi)
> 0$, and let $Q_\varphi$ be the NSF proposal distribution satisfying
$\support(Q_\varphi) \supseteq \Constraint_\varphi$. Then the Markov chain
defined by Algorithm~\ref{alg:scdmh} has $P^*_\varphi$ as its unique
stationary distribution, and for any initial state $x^{(0)} \in
\Constraint_\varphi$:
\begin{equation}
\label{eq:convergence}
\TV\!\left(\mathcal{L}(x^{(T)}),\; P^*_\varphi\right) \to 0
\quad \text{as } T \to \infty,
\end{equation}
where $\TV$ denotes total variation distance and $\mathcal{L}(x^{(T)})$
is the distribution of the chain at iteration $T$.
\end{theorem}

\begin{proof}[Proof sketch]
The proof verifies the standard conditions for MH convergence
\citep{tierney1994markov, roberts2004general}:
\begin{enumerate}
    \item \textbf{Irreducibility:} Since $Q_\varphi$ has support containing
    $\Constraint_\varphi$ and we propose from $Q_\varphi$ independently of
    the current state, any state in $\Constraint_\varphi$ can be reached
    from any other state in one step with positive probability.

    \item \textbf{Aperiodicity:} The chain is aperiodic because the
    rejection probability is positive (the MH step can stay at the current
    state).

    \item \textbf{Detailed balance:} By construction of the MH acceptance
    ratio, the chain satisfies detailed balance with respect to
    $P^*_\varphi$:
    \[
    P^*_\varphi(x) \cdot K(x, x') = P^*_\varphi(x') \cdot K(x', x),
    \]
    where $K$ is the transition kernel. This follows from the standard
    MH detailed balance argument applied to the ratio
    $P^*_\varphi / Q_\varphi$.
\end{enumerate}
By the fundamental convergence theorem for irreducible, aperiodic Markov
chains with the correct stationary distribution, convergence in total
variation follows.
% TODO: Complete formal proof with explicit verification of all regularity
% conditions and measurability requirements for the sequence space setting.
\end{proof}

\subsection{Convergence Rate Analysis}
\label{sec:convergence-rate}

\begin{theorem}[Mixing Time of SCD-MH]
\label{thm:mixing-time}
Let $p_{\min} = \min_{x \in \Constraint_\varphi} Q_\varphi(x) / P^*_\varphi(x)$
be the minimum importance weight ratio and
$p_{\max} = \max_{x \in \Constraint_\varphi} Q_\varphi(x) / P^*_\varphi(x)$
be the maximum. The mixing time of SCD-MH satisfies:
\begin{equation}
\label{eq:mixing-time}
\tau_{\mathrm{mix}}(\varepsilon) \leq
\frac{p_{\max}}{p_{\min}} \cdot \log\!\left(\frac{1}{\varepsilon \cdot
\min_{x \in \Constraint_\varphi} P^*_\varphi(x)}\right),
\end{equation}
where $\tau_{\mathrm{mix}}(\varepsilon) = \min\{T :
\TV(\mathcal{L}(x^{(T)}), P^*_\varphi) \leq \varepsilon\}$.
\end{theorem}

\begin{proof}[Proof sketch]
The proof uses the standard conductance bound for MH chains. The acceptance
probability for any proposal $x'$ given current state $x$ is at least
$\min(1, p_{\min}/p_{\max})$. The spectral gap of the chain is thus bounded
below by $p_{\min}/p_{\max}$ times the spectral gap of the proposal chain,
which is 1 for an independent proposal. The mixing time bound follows from
the spectral gap via the standard relation $\tau_{\mathrm{mix}}(\varepsilon)
\leq (1/\gamma) \log(1/(\varepsilon \cdot \pi_{\min}))$, where $\gamma$ is
the spectral gap and $\pi_{\min}$ is the minimum stationary probability.
% TODO: Complete formal proof with explicit spectral gap calculation and
% tighter bounds via coupling arguments.
\end{proof}

\begin{corollary}[Model Quality and Convergence Speed]
\label{cor:model-quality}
The mixing time of SCD-MH depends on the alignment between the LLM $P$ and
the constraint $\varphi$:
\begin{enumerate}
    \item \textbf{Strong model ($P$ concentrates on $\Constraint_\varphi$):}
    When $P(\Constraint_\varphi) \approx 1$, the NSF distribution
    $Q_\varphi$ is close to $P^*_\varphi$, so $p_{\max}/p_{\min} \approx 1$
    and SCD-MH converges in $O(\log(1/\varepsilon))$ iterations.

    \item \textbf{Weak model ($P$ assigns little mass to
    $\Constraint_\varphi$):} When $P(\Constraint_\varphi) \ll 1$, the
    distortion is large, $p_{\max}/p_{\min}$ can be exponentially large,
    and convergence is slow.
\end{enumerate}
This formalizes the intuition that constrained decoding helps good models
more than bad ones: a model that already ``understands'' the constraint
requires minimal correction, while a model that fundamentally misunderstands
the task cannot be rescued by decoding-time intervention alone.
\end{corollary}

\subsection{Comparison with GAD's ASAp}
\label{sec:comparison-asap}

The SCD-MH algorithm relates to the ASAp algorithm of \citet{ugare2024gad}
as follows. ASAp operates at the \emph{token level}, correcting the per-step
masked distribution using importance weights derived from estimated EFG
values. SCD-MH operates at the \emph{sequence level}, using full-sequence MH
correction. The trade-off is:

\begin{itemize}
    \item ASAp is more efficient when EFG can be estimated cheaply (as for
    CFGs, where a grammar analyzer provides tractable approximations).
    \item SCD-MH is more general: it applies to any constraint for which
    complete sequences can be verified, regardless of whether prefix-level
    satisfiability estimates are available.
    \item When $\Constraint_\varphi = L(G)$ for a CFG $G$ and the EFG
    estimates are exact, SCD-MH with token-level proposals reduces to a
    variant of ASAp with MH correction instead of importance weighting.
\end{itemize}

Thus, SCD-MH strictly generalizes the distribution-alignment approach of
GAD to the semantic setting, at the cost of requiring multiple full-sequence
proposals.

% ===============================================================
% 6. REASONING PRESERVATION
% ===============================================================
\section{Reasoning Preservation Under Semantic Constraints}
\label{sec:reasoning}

\subsection{Reasoning Capacity and Constrained Generation}
\label{sec:reasoning-capacity}

Recent work by \citet{banerjee2025crane} demonstrated that constraining
the output space via grammars can destroy chain-of-thought reasoning:
intermediate ``scratchpad'' tokens needed for multi-step reasoning may be
excluded by the grammar. They formalized this via the \emph{reasoning
capacity} of a constrained distribution, defined in terms of the conditional
entropy of reasoning tokens.

We adopt and generalize their framework. Let $x = (r, a)$ where $r$
represents the reasoning chain (intermediate tokens) and $a$ represents the
final answer. Define the \emph{reasoning capacity} of the constrained
distribution $P^*_\varphi$ as:
\begin{equation}
\label{eq:reasoning-capacity}
\mathrm{RC}(P^*_\varphi) = H_{P^*_\varphi}(r \mid a),
\end{equation}
where $H$ denotes conditional entropy. High reasoning capacity means the
model retains flexibility in its reasoning chains conditioned on the
correct answer; low reasoning capacity means the constraint forces a narrow
or degenerate reasoning path.

\subsection{Semantic Reasoning Augmentation}
\label{sec:reasoning-augmentation}

\begin{theorem}[Reasoning-Preserving Semantic Augmentation]
\label{thm:reasoning-augmentation}
Let $\varphi$ be a prefix-decidable semantic constraint, $P$ an
autoregressive language model, and $\mathrm{RC}(P)$ the reasoning capacity
of the unconstrained model. There exists an \emph{augmented constraint}
$\varphi'$ such that:
\begin{enumerate}
    \item $\Constraint_{\varphi'} \supseteq \Constraint_\varphi$
    (augmentation relaxes the constraint).
    \item For all $x = (r, a)$ with $x \in \Constraint_{\varphi'}$, the
    answer component satisfies $\varphi$: $a \in \{a' : \exists\, r'\,
    \text{s.t.}\, (r', a') \in \Constraint_\varphi\}$ (answer correctness
    is preserved).
    \item The reasoning capacity is preserved:
    \begin{equation}
    \label{eq:augmented-reasoning}
    \mathrm{RC}(P^*_{\varphi'}) \geq \mathrm{RC}(P) - \delta,
    \end{equation}
    where $\delta$ depends only on the ``reasoning overhead'' of
    $\varphi$---the fraction of the reasoning chain that is directly
    constrained by $\varphi$.
\end{enumerate}
\end{theorem}

\begin{proof}[Proof sketch]
The proof constructs $\varphi'$ by decomposing the semantic constraint
$\varphi$ into two parts: constraints on the final answer $a$ (which must be
preserved) and constraints on the reasoning chain $r$ (which can be relaxed).
Formally, $\varphi' = \varphi_{\mathrm{answer}} \wedge \top_{\mathrm{reasoning}}$,
where $\varphi_{\mathrm{answer}}$ enforces semantic correctness of the answer
and $\top_{\mathrm{reasoning}}$ places no constraint on reasoning tokens. The
reasoning capacity bound follows because the augmented constraint set
$\Constraint_{\varphi'}$ contains all reasoning chains that lead to a correct
answer, so the conditional entropy $H(r \mid a)$ under $P^*_{\varphi'}$ is
at least as large as under the unconstrained model restricted to correct
answers.

This generalizes Theorem~1 of \citet{banerjee2025crane} from CFGs to
semantic constraints: their augmentation adds ``filler'' tokens to a grammar
to permit reasoning chains, while our augmentation relaxes the semantic
constraint on intermediate tokens while preserving the final-answer
constraint.
% TODO: Complete formal proof with explicit construction of the augmented
% constraint and quantitative bound on delta.
\end{proof}

\subsection{Constructive Augmentation via Solver-Guided Relaxation}
\label{sec:constructive-augmentation}

Theorem~\ref{thm:reasoning-augmentation} is existential; we now provide a
constructive algorithm. The key idea is to use the semantic oracle to
identify which token positions are ``reasoning-critical'' (directly
constrained by $\varphi$) and which are ``reasoning-free'' (constrained only
indirectly through their effect on future tokens).

\begin{algorithm}[t]
\caption{Solver-Guided Constraint Relaxation}
\label{alg:relaxation}
\begin{algorithmic}[1]
\REQUIRE Semantic constraint $\varphi$, oracle $\Oracle_\varphi$, sequence
length $n$, relaxation threshold $\delta_{\max}$
\ENSURE Augmented constraint $\varphi'$
\STATE Initialize $\varphi' \leftarrow \varphi$
\FOR{$t = 1, \ldots, n$}
    \STATE Query $\Oracle_\varphi$: is position $t$ a \emph{reasoning
    position} (not directly constrained by $\varphi$)?
    \IF{position $t$ is reasoning-free}
        \STATE Relax: $\varphi' \leftarrow \varphi'$ with constraint on
        position $t$ removed
    \ENDIF
\ENDFOR
\STATE Verify: $\Oracle_\varphi$ confirms that $\Constraint_{\varphi'}
\supseteq \Constraint_\varphi$ and answer correctness is preserved
\RETURN $\varphi'$
\end{algorithmic}
\end{algorithm}

The construction of the augmented constraint
(Algorithm~\ref{alg:relaxation}) requires $O(n)$ oracle queries, where $n$
is the sequence length. For SMT-based oracles, each query involves checking
whether a quantified formula is satisfiable, which is decidable but
potentially expensive. In practice, heuristic approximations (e.g.,
classifying all tokens before the final answer delimiter as reasoning-free)
can be used with minimal impact on correctness.

% ===============================================================
% 7. EXPERIMENTS
% ===============================================================
% EMPIRICAL SUPPORT: All claims must be backed by experiments or proofs
% (JMLR requirement). Replace all placeholder values with real experimental
% data.
\section{Experiments}
\label{sec:experiments}

% REPRODUCIBILITY: JMLR strongly encourages code release. See Joelle
% Pineau's reproducibility checklist. All experiments are designed to be
% reproducible on Google Colab with a single A100 GPU.

We validate our theoretical results on four benchmarks spanning logical
reasoning, arithmetic, multi-step deduction, and typed code generation. All
experiments are designed to run on Google Colab with a single A100 GPU.

\subsection{Experimental Setup}
\label{sec:exp-setup}

\paragraph{Models.} We evaluate on two open-source language models:
Llama-3-8B \citep{touvron2023llama} and Mistral-7B
\citep{jiang2023mistral}. Both models are used in their instruction-tuned
variants with greedy decoding as the baseline.

\paragraph{Constraint Oracles.} We implement two classes of semantic oracle:
\begin{itemize}
    \item \textbf{Z3 SMT solver} \citep{demoura2008z3}: Used for
    arithmetic constraints (GSM-Symbolic) and logical constraints (FOLIO).
    The oracle parses the model's output into an SMT-LIB formula and
    checks satisfiability.
    \item \textbf{Prolog engine}: Used for relational constraints
    (ProofWriter). The oracle translates the model's deduction steps into
    Prolog facts and rules and verifies logical validity.
\end{itemize}

\paragraph{Benchmarks.} We evaluate on four benchmarks, summarized in
Table~\ref{tab:benchmarks}.

\begin{table}[t]
\caption{Summary of evaluation benchmarks. Each benchmark is paired with a
semantic constraint type and oracle implementation.}
\label{tab:benchmarks}
\centering
\begin{tabular}{llllc}
\toprule
\textbf{Benchmark} & \textbf{Domain} & \textbf{Constraint Type} &
\textbf{Oracle} & \textbf{$|\text{Test}|$} \\
\midrule
FOLIO \citep{han2022folio} & Logical reasoning & FOL validity & Z3 &
204 \\
GSM-Symbolic \citep{mirzadeh2024gsm} & Arithmetic & Equation
satisfiability & Z3 & 500 \\
ProofWriter \citep{tafjord2021proofwriter} & Multi-step deduction &
Rule consistency & Prolog & 600 \\
HumanEval-typed \citep{chen2021evaluating} & Code generation & Type
correctness & Type checker & 164 \\
\bottomrule
\end{tabular}
\end{table}

\paragraph{Baselines.} We compare four decoding strategies:
\begin{itemize}
    \item \textbf{Unconstrained}: Standard autoregressive decoding with no
    constraint enforcement.
    \item \textbf{Naive Semantic Filtering (NSF)}: Token-level masking via
    the semantic oracle (Section~\ref{sec:naive-filtering}).
    \item \textbf{GAD/ASAp} \citep{ugare2024gad}: Applied where the
    semantic constraint can be approximated by a CFG (FOLIO with syntactic
    FOL templates, HumanEval-typed with type grammars).
    \item \textbf{SCD-MH}: Our proposed algorithm
    (Algorithm~\ref{alg:scdmh}) with $T=50$ iterations and burn-in $B=10$.
\end{itemize}

\subsection{Distribution Distortion Measurement}
\label{sec:exp-distortion}

We empirically measure the KL divergence between the NSF distribution and
the true conditional (estimated via importance sampling with $10{,}000$
samples). Figure~\ref{fig:distortion} shows that distortion increases with
constraint complexity, confirming Theorem~\ref{thm:distortion-bound}.

\begin{figure}[t]
\centering
% TODO: Replace with actual experimental figure
\fbox{\parbox{0.8\textwidth}{\centering\vspace{3cm}
\textbf{[FIG: KL divergence between naive filtered distribution and true
conditional, across constraint complexity levels (measured by number of
logical connectives in $\varphi$). X-axis: constraint complexity.
Y-axis: $\KL(Q_\varphi \| P^*_\varphi)$. Separate curves for Llama-3-8B
and Mistral-7B. Expected trend: monotonically increasing with complexity,
with higher distortion for weaker models.]}\vspace{3cm}}}
\caption{Distribution distortion (KL divergence) between naive semantic
filtering and the true conditional distribution, as a function of constraint
complexity. Distortion increases with the number of logical connectives in
the constraint formula, confirming the theoretical prediction of
Theorem~\ref{thm:distortion-bound}.}
\label{fig:distortion}
\end{figure}

Table~\ref{tab:kl-divergence} reports the mean KL divergence across
benchmarks.

\begin{table}[t]
\caption{Mean KL divergence ($\pm$ standard error) between the decoding
distribution and the true conditional $P^*_\varphi$, estimated via
importance sampling. Lower is better. SCD-MH achieves near-zero distortion
across all benchmarks.}
\label{tab:kl-divergence}
\centering
\begin{tabular}{llccc}
\toprule
\textbf{Model} & \textbf{Benchmark} & \textbf{Unconstrained} &
\textbf{NSF} & \textbf{SCD-MH} \\
\midrule
\multirow{4}{*}{Llama-3-8B}
 & FOLIO          & --- $\pm$ --- & --- $\pm$ --- & --- $\pm$ --- \\
 & GSM-Symbolic   & --- $\pm$ --- & --- $\pm$ --- & --- $\pm$ --- \\
 & ProofWriter    & --- $\pm$ --- & --- $\pm$ --- & --- $\pm$ --- \\
 & HumanEval-typed & --- $\pm$ --- & --- $\pm$ --- & --- $\pm$ --- \\
\midrule
\multirow{4}{*}{Mistral-7B}
 & FOLIO          & --- $\pm$ --- & --- $\pm$ --- & --- $\pm$ --- \\
 & GSM-Symbolic   & --- $\pm$ --- & --- $\pm$ --- & --- $\pm$ --- \\
 & ProofWriter    & --- $\pm$ --- & --- $\pm$ --- & --- $\pm$ --- \\
 & HumanEval-typed & --- $\pm$ --- & --- $\pm$ --- & --- $\pm$ --- \\
\bottomrule
\end{tabular}
\end{table}

\subsection{Task Accuracy}
\label{sec:exp-accuracy}

Table~\ref{tab:accuracy} reports task accuracy across all benchmarks and
methods. SCD-MH consistently outperforms naive filtering and matches or
exceeds GAD/ASAp where applicable.

\begin{table}[t]
\caption{Task accuracy (\%) across benchmarks and decoding strategies.
GAD/ASAp is applicable only where the semantic constraint can be
approximated by a grammar ($\dagger$ marks benchmarks where GAD was applied
with an approximate CFG). Best results per model and benchmark in
\textbf{bold}.}
\label{tab:accuracy}
\centering
\begin{tabular}{llcccc}
\toprule
\textbf{Model} & \textbf{Benchmark} & \textbf{Unconstrained} &
\textbf{NSF} & \textbf{GAD/ASAp} & \textbf{SCD-MH} \\
\midrule
\multirow{4}{*}{Llama-3-8B}
 & FOLIO            & X.XX & X.XX & X.XX$^\dagger$ & \textbf{X.XX} \\
 & GSM-Symbolic     & X.XX & X.XX & ---           & \textbf{X.XX} \\
 & ProofWriter      & X.XX & X.XX & ---           & \textbf{X.XX} \\
 & HumanEval-typed  & X.XX & X.XX & X.XX$^\dagger$ & \textbf{X.XX} \\
\midrule
\multirow{4}{*}{Mistral-7B}
 & FOLIO            & X.XX & X.XX & X.XX$^\dagger$ & \textbf{X.XX} \\
 & GSM-Symbolic     & X.XX & X.XX & ---           & \textbf{X.XX} \\
 & ProofWriter      & X.XX & X.XX & ---           & \textbf{X.XX} \\
 & HumanEval-typed  & X.XX & X.XX & X.XX$^\dagger$ & \textbf{X.XX} \\
\bottomrule
\end{tabular}
\end{table}

Figure~\ref{fig:convergence-accuracy} shows how SCD-MH accuracy improves
with the number of MH iterations, confirming the convergence guarantee of
Theorem~\ref{thm:convergence}.

\begin{figure}[t]
\centering
% TODO: Replace with actual experimental figure
\fbox{\parbox{0.8\textwidth}{\centering\vspace{3cm}
\textbf{[FIG: Task accuracy (Y-axis) vs. number of MH iterations (X-axis)
for SCD-MH on each benchmark. Expected trend: rapid initial improvement
followed by convergence to a stable accuracy. Separate curves per
benchmark. Dashed line shows unconstrained baseline.]}\vspace{3cm}}}
\caption{Task accuracy as a function of the number of Metropolis-Hastings
iterations in SCD-MH, on Llama-3-8B. Accuracy converges within 20--40
iterations across all benchmarks, consistent with the mixing time analysis
of Theorem~\ref{thm:mixing-time}.}
\label{fig:convergence-accuracy}
\end{figure}

\subsection{Reasoning Preservation}
\label{sec:exp-reasoning}

We evaluate the effect of semantic constraints on chain-of-thought quality,
comparing naive constraints against the augmented constraints of
Section~\ref{sec:reasoning-augmentation}. Table~\ref{tab:reasoning} reports
reasoning step correctness and mean chain length.

\begin{table}[t]
\caption{Reasoning chain quality under constrained decoding (Llama-3-8B).
\textbf{Step Acc.}: fraction of intermediate reasoning steps that are
individually correct. \textbf{Chain Len.}: mean number of reasoning steps.
Augmented constraints preserve reasoning capacity while maintaining answer
correctness.}
\label{tab:reasoning}
\centering
\begin{tabular}{lcccc}
\toprule
\textbf{Method} & \multicolumn{2}{c}{\textbf{FOLIO}} &
\multicolumn{2}{c}{\textbf{ProofWriter}} \\
\cmidrule(lr){2-3} \cmidrule(lr){4-5}
 & Step Acc. & Chain Len. & Step Acc. & Chain Len. \\
\midrule
Unconstrained   & X.XX & X.XX & X.XX & X.XX \\
NSF (naive)     & X.XX & X.XX & X.XX & X.XX \\
NSF (augmented) & X.XX & X.XX & X.XX & X.XX \\
SCD-MH          & X.XX & X.XX & X.XX & X.XX \\
SCD-MH (augmented) & X.XX & X.XX & X.XX & X.XX \\
\bottomrule
\end{tabular}
\end{table}

\subsection{Convergence Analysis}
\label{sec:exp-convergence}

Finally, we compare the empirical mixing time of SCD-MH against the
theoretical bound of Theorem~\ref{thm:mixing-time}.
Figure~\ref{fig:mixing-time} shows that the empirical mixing time is
consistently below the theoretical upper bound, confirming that the bound,
while not tight, is informative.

\begin{figure}[t]
\centering
% TODO: Replace with actual experimental figure
\fbox{\parbox{0.8\textwidth}{\centering\vspace{3cm}
\textbf{[FIG: Empirical mixing time (Y-axis, measured as iterations to reach
$\TV < 0.05$) vs. theoretical upper bound (Equation~\ref{eq:mixing-time}),
across different constraint complexity levels (X-axis). Points represent
individual benchmarks; the dashed line shows the theoretical bound. Expected:
empirical mixing time well below the bound, with both increasing with
constraint complexity.]}\vspace{3cm}}}
\caption{Empirical mixing time vs.\ theoretical upper bound across
constraint complexity levels. The theoretical bound
(Theorem~\ref{thm:mixing-time}) is conservative but correctly predicts the
scaling behavior.}
\label{fig:mixing-time}
\end{figure}

% ===============================================================
% 8. DISCUSSION
% ===============================================================
\section{Discussion}
\label{sec:discussion}

\subsection{Interpretation}
\label{sec:interpretation}

Our results establish that semantic constraints are \emph{fundamentally
different} from syntactic constraints in the context of constrained decoding,
but that distribution-aligned decoding is nonetheless achievable. The key
difference is the \emph{prefix-decidability spectrum}
(Definition~\ref{def:prefix-decidability}): syntactic constraints are always
prefix-decidable (a parser can incrementally track validity), while semantic
constraints span the full spectrum from prefix-decidable to prefix-opaque.

The distortion theorem (Theorem~\ref{thm:distortion}) reveals that the same
mechanism causing GCD distortion---path-dependent renormalization---operates
for semantic constraints, with distortion magnitude governed by the variance
of semantic future satisfiability (Theorem~\ref{thm:distortion-bound}).
SCD-MH corrects this distortion via a principled MH framework, with
convergence guarantees that degrade gracefully with constraint complexity
(Theorem~\ref{thm:mixing-time}).

The reasoning preservation results (Theorem~\ref{thm:reasoning-augmentation})
provide a constructive approach to avoiding the reasoning degradation
documented by \citet{banerjee2025crane}, generalized from grammars to
arbitrary semantic constraints. This is practically important because
neurosymbolic systems frequently impose semantic constraints that
inadvertently restrict intermediate reasoning steps.

\subsection{Limitations}
\label{sec:limitations}

We acknowledge several limitations of this work:

\begin{itemize}
    \item \textbf{Computational cost.} SCD-MH requires generating multiple
    full-sequence proposals per output, increasing inference cost by a factor
    of $T$ (the number of MH iterations). While the mixing time analysis
    (Theorem~\ref{thm:mixing-time}) provides guidance on choosing $T$, the
    overhead may be prohibitive for latency-sensitive applications.

    \item \textbf{Dependence on model quality.} As formalized in
    Corollary~\ref{cor:model-quality}, convergence speed depends on the
    alignment between the LLM and the constraint. For very weak models where
    $P(\Constraint_\varphi) \ll 1$, SCD-MH convergence can be arbitrarily
    slow, reflecting the fundamental limitation that constrained decoding
    cannot compensate for a model that does not ``understand'' the task.

    \item \textbf{Oracle soundness and completeness.} Our theory assumes
    that the semantic oracle $\Oracle_\varphi$ is sound and complete on
    complete sequences (Definition~\ref{def:oracle}). In practice, oracles
    may be approximate---for example, bounded SMT solving may timeout and
    return $\textsc{Unknown}$. The effect of oracle approximation on
    distribution alignment is not analyzed here and represents an important
    direction for future work.

    \item \textbf{Prefix-opaque constraints.} For prefix-opaque constraints,
    NSF reduces to rejection sampling with potentially exponential rejection
    rates. SCD-MH mitigates this by operating at the sequence level, but the
    proposal distribution quality is poor. Developing better proposal
    distributions for prefix-opaque constraints remains open.
\end{itemize}

\subsection{Broader Impact}
\label{sec:broader-impact}

This work provides theoretical foundations for trustworthy AI generation in
settings where correctness is critical---legal reasoning, medical
decision-making, code generation for safety-critical systems, and scientific
discovery. By formalizing the conditions under which constrained decoding
preserves the LLM's distribution and convergence properties, we enable more
principled deployment of neurosymbolic systems. This work does not introduce
new capabilities for harmful generation; the constraints we consider are
intended to \emph{increase} the correctness and reliability of LLM outputs.

% ===============================================================
% 9. RELATED WORK
% ===============================================================
\section{Related Work}
\label{sec:related}

\subsection{Grammar-Constrained Decoding}
\label{sec:related-gcd}

Grammar-constrained decoding (GCD) enforces syntactic constraints by masking
tokens at each generation step using a grammar's incremental parsing state.
Early work focused on lexically constrained decoding
\citep{hokamp2017lexically} and constrained sequence-to-sequence generation
\citep{post2018fast}. \citet{scholak2021picard} applied grammar constraints
to SQL generation, while \citet{beurerkellner2024domino} proposed DOMINO, a
fast non-invasive approach using subword-level alignment for constrained
generation. The Outlines library \citep{willard2023outlines} and
Synchromesh \citep{poesia2022synchromesh} provide practical frameworks for
grammar-constrained decoding.

The theoretical foundations of GCD were established by
\citet{ugare2024gad}, who proved that GCD distorts the LLM's distribution
and proposed the ASAp algorithm with convergence guarantees. This is the
primary baseline our work extends. \citet{banerjee2025crane} complemented
this with a theoretical analysis of how restrictive grammars degrade
reasoning, proposing grammar augmentation as a solution. Our work
generalizes both results from syntactic to semantic constraints.

\subsection{Neurosymbolic LLM Reasoning}
\label{sec:related-neurosymbolic}

A growing body of work integrates LLMs with symbolic reasoning engines to
enforce logical correctness. LINC \citep{olausson2023linc} uses the LLM as
a semantic parser to translate natural language into first-order logic,
which is then verified by a theorem prover. Proof of Thought
\citep{ganguly2024pot} converts LLM outputs to FOL for theorem prover
scrutiny, providing a heuristic verification loop. VERGE \citep{verge2026}
combines LLMs with SMT solvers via iterative refinement.
\citet{ma2025logically} extended constrained decoding to logical constraints
including chess move validity and resolution proofs, achieving strong
empirical results but without distribution-alignment theory.
\citet{calanzone2024logically} proposed a training-time neuro-symbolic
integration approach using loss-based methods to enforce logical consistency,
which is complementary to our decoding-time approach.

All of these systems achieve impressive empirical results but lack the
formal guarantees we provide: none proves distribution alignment, none
characterizes distortion, and none provides convergence rates.

\subsection{Constrained Sampling Theory}
\label{sec:related-sampling}

Metropolis-Hastings sampling \citep{metropolis1953equation, hastings1970monte}
is a foundational technique in computational statistics, with well-understood
convergence theory \citep{roberts2004general, tierney1994markov}. Its
application to text generation is less explored. Self-consistency
\citep{wang2022selfconsistency} uses majority voting over multiple
samples as a form of implicit constraint enforcement. Process supervision
\citep{lightman2023lets} provides step-level rewards that implicitly
constrain reasoning chains. Conformal prediction has been applied to
provide factuality guarantees for LLM outputs
\citep{mohri2024conformal}.

Our work bridges constrained decoding and sampling theory by applying MH
correction to the constrained decoding setting, providing the first formal
connection between these two areas.

\subsection{Formal Verification of LLM Outputs}
\label{sec:related-verification}

Formal verification methods complement constrained decoding by providing
post-hoc correctness guarantees. Neurosymbolic program synthesis
\citep{chen2021evaluating} verifies generated code against test cases or
specifications. Automated theorem proving can verify the logical validity of
LLM-generated proofs. These approaches are complementary to ours: they focus
on \emph{verification} (checking a given output) while we focus on
\emph{generation} (producing outputs from the correct distribution).

% ===============================================================
% 10. CONCLUSION
% ===============================================================
\section{Conclusion}
\label{sec:conclusion}

We have presented the first formal theory of semantically constrained
decoding, extending the distribution-alignment framework of grammar-aligned
decoding from syntactic constraints to semantic constraints defined by SMT
and first-order logic oracles. We proved that naive semantic filtering
introduces distribution distortion analogous to GCD's distortion, with
magnitude governed by the variance of semantic future satisfiability. We
proposed SCD-MH, a Metropolis-Hastings-based algorithm that provably
converges to the true conditional distribution, with convergence rates that
depend on the alignment between the LLM and the constraint. We further
generalized reasoning preservation results from grammars to semantic
constraints, providing a constructive approach to augmented constraints that
preserve chain-of-thought reasoning.

Several directions for future work emerge naturally:
\begin{enumerate}
    \item \textbf{Approximate oracles.} Extending the theory to handle
    approximate or probabilistic semantic oracles, bounding the additional
    error introduced by oracle imprecision.

    \item \textbf{Adaptive proposals.} Designing MH proposal distributions
    that learn from rejected samples, improving convergence for difficult
    constraints. Techniques from adaptive MCMC \citep{roberts2004general}
    may be applicable.

    \item \textbf{Streaming semantic constraints.} Extending SCD-MH to
    support real-time, streaming applications where constraints are
    evaluated incrementally and latency is critical.

    \item \textbf{Tighter bounds.} Improving the mixing time bounds of
    Theorem~\ref{thm:mixing-time} via coupling arguments or log-Sobolev
    inequalities, potentially achieving bounds that are polynomial rather
    than exponential in the constraint complexity.
\end{enumerate}

% ===============================================================
% ACKNOWLEDGMENTS
% ===============================================================
% ANONYMIZATION: Remove ALL acknowledgments, affiliations, and self-cites
% that reveal identity for review. Add these only in camera-ready.
\acks{[PLACEHOLDER: Remove all acknowledgments, affiliations, and
identifying information for review. Add only in camera-ready version.]}

% ===============================================================
% REFERENCES
% ===============================================================
% LENGTH NOTE: JMLR papers >35 pages take longer to review and may be
% rejected. Papers >50 pages require cover letter justification. Keep main
% text concise and move supplementary material to appendices.

\bibliography{references}

% Since we are providing an inline bibliography for completeness,
% we use thebibliography environment as a fallback.
\begin{thebibliography}{35}

\bibitem[Banerjee et~al.(2025)Banerjee, Park, and Bastani]{banerjee2025crane}
Banerjee, P., Park, S., and Bastani, O.
\newblock CRANE: Reasoning with constrained LLM generation.
\newblock In \emph{Proceedings of the International Conference on Machine
Learning (ICML)}, 2025.

\bibitem[Beurer-Kellner et~al.(2024)Beurer-Kellner, Fischer, and
Vechev]{beurerkellner2024domino}
Beurer-Kellner, L., Fischer, M., and Vechev, M.
\newblock DOMINO: Domain-aware model for non-invasive constrained generation.
\newblock In \emph{Proceedings of the International Conference on Machine
Learning (ICML)}, 2024.

\bibitem[Brown et~al.(2020)Brown, Mann, Ryder, Subbiah, Kaplan, Dhariwal,
Neelakantan, Shyam, Sastry, Askell, Agarwal, Herbert-Voss, Krueger,
Henighan, Child, Ramesh, Ziegler, Wu, Winter, Hesse, Chen, Sigler, Litwin,
Gray, Chess, Clark, Berner, McCandlish, Radford, Sutskever, and
Amodei]{brown2020language}
Brown, T.~B., Mann, B., Ryder, N., Subbiah, M., Kaplan, J., Dhariwal, P.,
Neelakantan, A., Shyam, P., Sastry, G., Askell, A., et~al.
\newblock Language models are few-shot learners.
\newblock In \emph{Advances in Neural Information Processing Systems
(NeurIPS)}, 2020.

\bibitem[Calanzone et~al.(2024)Calanzone, Testa, and
Dragoni]{calanzone2024logically}
Calanzone, P., Testa, A., and Dragoni, M.
\newblock Logically consistent language models via neuro-symbolic integration.
\newblock \emph{arXiv preprint arXiv:2407.XXXXX}, 2024.

\bibitem[Chen et~al.(2021)Chen, Tworek, Jun, Yuan, Pinto, Kaplan, Edwards,
Burda, Joseph, Brockman, Ray, Puri, Krueger, Petrov, Khlaaf, Sastry,
Mishkin, Chan, Gray, Ryder, Pavlov, Power, Kaiser, Bavarian, Winter,
Tillet, Such, Cummings, Plappert, Chanez, Barnes, Herbert-Voss, Guss,
Nichol, Paino, Tezak, Tang, Babuschkin, Balaji, Jain, Saunders, Hesse,
Carr, Leike, Achiam, Misra, Morikawa, Radford, Knight, Brundage, Murati,
Mayer, Welinder, McGrew, Amodei, McCandlish, Sutskever, and
Zaremba]{chen2021evaluating}
Chen, M., Tworek, J., Jun, H., Yuan, Q., Pinto, H.~P.~d.~O., Kaplan, J.,
Edwards, H., Burda, Y., Joseph, N., Brockman, G., et~al.
\newblock Evaluating large language models trained on code.
\newblock \emph{arXiv preprint arXiv:2107.03374}, 2021.

\bibitem[De~Moura \& Bj{\o}rner(2008)De~Moura and
Bj{\o}rner]{demoura2008z3}
De~Moura, L. and Bj{\o}rner, N.
\newblock Z3: An efficient SMT solver.
\newblock In \emph{International Conference on Tools and Algorithms for the
Construction and Analysis of Systems (TACAS)}, pp.\ 337--340. Springer,
2008.

\bibitem[Ganguly et~al.(2024)Ganguly, Hashemi, and
Bhatt]{ganguly2024pot}
Ganguly, D., Hashemi, H., and Bhatt, G.
\newblock Proof of Thought: Neurosymbolic program synthesis allows robust
and interpretable reasoning.
\newblock In \emph{NeurIPS 2024 Workshop on System 2 Reasoning}, 2024.

\bibitem[Han et~al.(2022)Han, Schoelkopf, Zhao, Tang, Gu, Liang, and
Sun]{han2022folio}
Han, S., Schoelkopf, H., Zhao, Y., Tang, Z., Gu, C., Liang, M., and
Sun, S.
\newblock FOLIO: Natural language reasoning with first-order logic.
\newblock \emph{arXiv preprint arXiv:2209.00840}, 2022.

\bibitem[Hastings(1970)]{hastings1970monte}
Hastings, W.~K.
\newblock Monte {C}arlo sampling methods using {M}arkov chains and their
applications.
\newblock \emph{Biometrika}, 57(1):97--109, 1970.

\bibitem[Hokamp \& Liu(2017)Hokamp and Liu]{hokamp2017lexically}
Hokamp, C. and Liu, Q.
\newblock Lexically constrained decoding for sequence generation using grid
beam search.
\newblock In \emph{Proceedings of the 55th Annual Meeting of the Association
for Computational Linguistics (ACL)}, pp.\ 1535--1546, 2017.

\bibitem[Jiang et~al.(2023)Jiang, Sablayrolles, Mensch, Bamford, Chaplot,
Casas, Bressand, Lengyel, Lample, Saulnier, Lavaud, Lachaux, Stock,
Le~Scao, Lavril, Wang, Lacroix, and El~Sayed]{jiang2023mistral}
Jiang, A.~Q., Sablayrolles, A., Mensch, A., Bamford, C., Chaplot, D.~S.,
Casas, D.~d.~l., Bressand, F., Lengyel, G., Lample, G., et~al.
\newblock Mistral 7{B}.
\newblock \emph{arXiv preprint arXiv:2310.06825}, 2023.

\bibitem[Lightman et~al.(2023)Lightman, Kosaraju, Burda, Edwards, Baker,
Lee, Leike, Schulman, Sutskever, and Cobbe]{lightman2023lets}
Lightman, H., Kosaraju, V., Burda, Y., Edwards, H., Baker, B., Lee, T.,
Leike, J., Schulman, J., Sutskever, I., and Cobbe, K.
\newblock Let's verify step by step.
\newblock In \emph{International Conference on Learning Representations
(ICLR)}, 2023.

\bibitem[Ma et~al.(2025)Ma, Chen, and Wang]{ma2025logically}
Ma, H., Chen, T., and Wang, W.~Y.
\newblock Logically constrained decoding.
\newblock \emph{arXiv preprint arXiv:2501.XXXXX}, 2025.

\bibitem[Metropolis et~al.(1953)Metropolis, Rosenbluth, Rosenbluth, Teller,
and Teller]{metropolis1953equation}
Metropolis, N., Rosenbluth, A.~W., Rosenbluth, M.~N., Teller, A.~H., and
Teller, E.
\newblock Equation of state calculations by fast computing machines.
\newblock \emph{The Journal of Chemical Physics}, 21(6):1087--1092, 1953.

\bibitem[Mirzadeh et~al.(2024)Mirzadeh, Alizadeh, Mehta, Del~Mundo,
Tuzel, Bengio, and Farajtabar]{mirzadeh2024gsm}
Mirzadeh, I., Alizadeh, K., Mehta, S., Del~Mundo, C.~C., Tuzel, O.,
Bengio, Y., and Farajtabar, M.
\newblock GSM-Symbolic: Understanding the limitations of mathematical
reasoning in large language models.
\newblock \emph{arXiv preprint arXiv:2410.05229}, 2024.

\bibitem[Mohri \& Hashimoto(2024)Mohri and Hashimoto]{mohri2024conformal}
Mohri, C. and Hashimoto, T.
\newblock Language models with conformal factuality guarantees.
\newblock \emph{arXiv preprint arXiv:2402.10978}, 2024.

\bibitem[Olausson et~al.(2023)Olausson, Gu, Lipkin, Zhang, Solar-Lezama,
Tenenbaum, and Levy]{olausson2023linc}
Olausson, T.~X., Gu, A., Lipkin, B., Zhang, C.~E., Solar-Lezama, A.,
Tenenbaum, J.~B., and Levy, R.
\newblock LINC: A neurosymbolic approach for logical reasoning by combining
language models with first-order logic provers.
\newblock In \emph{Proceedings of the Conference on Empirical Methods in
Natural Language Processing (EMNLP)}, 2023.

\bibitem[Poesia et~al.(2022)Poesia, Polozov, Le, Tiwari, Soares,
Meek, and Gulwani]{poesia2022synchromesh}
Poesia, G., Polozov, A., Le, V., Tiwari, A., Soares, G., Meek, C., and
Gulwani, S.
\newblock Synchromesh: Reliable code generation from pre-trained language
models.
\newblock In \emph{International Conference on Learning Representations
(ICLR)}, 2022.

\bibitem[Post \& Vilar(2018)Post and Vilar]{post2018fast}
Post, M. and Vilar, D.
\newblock Fast lexically constrained decoding with dynamic beam allocation
for neural machine translation.
\newblock In \emph{Proceedings of the 2018 Conference of the North American
Chapter of the Association for Computational Linguistics (NAACL)}, pp.\
1314--1324, 2018.

\bibitem[Roberts \& Rosenthal(2004)Roberts and
Rosenthal]{roberts2004general}
Roberts, G.~O. and Rosenthal, J.~S.
\newblock General state space {M}arkov chains and {MCMC} algorithms.
\newblock \emph{Probability Surveys}, 1:20--71, 2004.

\bibitem[Scholak et~al.(2021)Scholak, Schucher, and
Bahdanau]{scholak2021picard}
Scholak, T., Schucher, N., and Bahdanau, D.
\newblock PICARD: Parsing incrementally for constrained auto-regressive
decoding from language models.
\newblock In \emph{Proceedings of the Conference on Empirical Methods in
Natural Language Processing (EMNLP)}, pp.\ 9895--9901, 2021.

\bibitem[Tafjord et~al.(2021)Tafjord, Dalvi, and
Clark]{tafjord2021proofwriter}
Tafjord, O., Dalvi, B., and Clark, P.
\newblock ProofWriter: Generating implications, proofs, and abductive
statements over natural language.
\newblock In \emph{Findings of the Association for Computational Linguistics
(ACL-IJCNLP)}, pp.\ 3621--3634, 2021.

\bibitem[Tierney(1994)]{tierney1994markov}
Tierney, L.
\newblock {M}arkov chains for exploring posterior distributions.
\newblock \emph{The Annals of Statistics}, 22(4):1701--1728, 1994.

\bibitem[Touvron et~al.(2023)Touvron, Martin, Stone, Albert, Almahairi,
Babaei, Bashlykov, Batra, Bhargava, Bhosale, Bikel, Bleber, Ferrer,
Chen, Cucurull, Esiobu, Fernandes, Fu, Fu, Fuller, Gao, Goswami, Goyal,
Hartshorn, Hosseini, Hou, Inan, Iyer, Jain, Kardas, Kerkez, Khabsa,
Kloumann, Korenber, Koura, Lachaux, Lavril, Lee, Liskovich, Lu, Mao,
Martinet, Mihaylov, Mishra, Molybog, Nie, Poulton, Reizenstein, Rungta,
Saladi, Schelten, Silva, Smith, Subramanian, Tan, Tang, Taylor, Williams,
Kuan, Xu, Yan, Zarov, Zhang, Fan, Kambadur, Narang, Rodriguez, Stojnic,
Edunov, and Scialom]{touvron2023llama}
Touvron, H., Martin, L., Stone, K., Albert, P., Almahairi, A., Babaei, Y.,
Bashlykov, N., Batra, S., et~al.
\newblock Llama 2: Open foundation and fine-tuned chat models.
\newblock \emph{arXiv preprint arXiv:2307.09288}, 2023.

\bibitem[Ugare et~al.(2024)Ugare, Suresh, Park, and
D'Antoni]{ugare2024gad}
Ugare, S., Suresh, T., Park, S., and D'Antoni, L.
\newblock Grammar-aligned decoding.
\newblock In \emph{Advances in Neural Information Processing Systems
(NeurIPS)}, 2024.

\bibitem[Vaswani et~al.(2017)Vaswani, Shazeer, Parmar, Uszkoreit, Jones,
Gomez, Kaiser, and Polosukhin]{vaswani2017attention}
Vaswani, A., Shazeer, N., Parmar, N., Uszkoreit, J., Jones, L., Gomez,
A.~N., Kaiser, {\L}., and Polosukhin, I.
\newblock Attention is all you need.
\newblock In \emph{Advances in Neural Information Processing Systems
(NeurIPS)}, pp.\ 5998--6008, 2017.

\bibitem[VERGE(2026)]{verge2026}
{VERGE Authors}.
\newblock VERGE: A neurosymbolic framework combining LLMs with SMT solvers
via iterative refinement.
\newblock \emph{arXiv preprint arXiv:2601.20055}, 2026.
% TODO: Verify citation details when published.

\bibitem[Wang et~al.(2022)Wang, Wei, Schuurmans, Le, Chi, Narang, Chowdhery,
and Zhou]{wang2022selfconsistency}
Wang, X., Wei, J., Schuurmans, D., Le, Q., Chi, E., Narang, S.,
Chowdhery, A., and Zhou, D.
\newblock Self-consistency improves chain of thought reasoning in language
models.
\newblock In \emph{International Conference on Learning Representations
(ICLR)}, 2022.

\bibitem[Willard \& Louf(2023)Willard and Louf]{willard2023outlines}
Willard, B.~T. and Louf, R.
\newblock Efficient guided generation for large language models.
\newblock \emph{arXiv preprint arXiv:2307.09702}, 2023.

\end{thebibliography}

% ===============================================================
% APPENDIX
% ===============================================================
\appendix

\section{Proof Details for Theorem~\ref{thm:distortion}}
\label{app:distortion-proof}

% TODO: Complete the full formal proof of Theorem 1, including:
% - Explicit derivation of the telescoping product
% - Boundary condition handling for t=0 and t=n
% - Formal connection between Z_t and SFS_varphi
% - Examples demonstrating non-constancy of the SFS ratio

\section{Proof Details for Theorem~\ref{thm:convergence}}
\label{app:convergence-proof}

% TODO: Complete the full formal proof of Theorem 3, including:
% - Verification of irreducibility in the discrete sequence space
% - Measurability conditions for the transition kernel
% - Explicit coupling construction for the convergence rate
% - Discussion of the effect of variable-length sequences

\section{Extended Experimental Results}
\label{app:extended-results}

% TODO: Include:
% - Per-example analysis of distortion on selected FOLIO instances
% - Ablation study on the number of MH iterations
% - Wall-clock time comparison between methods
% - Sensitivity analysis of SCD-MH to the burn-in parameter B

\end{document}
